% ======================================================================
% 00_preamble.tex: packages and formatting shared by the whole report.
% Built with XeLaTeX (fontspec reads UTF-8 and Unicode fonts directly).
% Names, title and dates come from info.txt, not from this file.
%
% WHAT YOU MAY WANT TO CHANGE (find the section by its banner below)
%   Font ........... the \setmainfont lines (any font installed on your computer).
%   Margins ........ the numbers in \usepackage[...]{geometry} (top, bottom, left, right).
%   Text size ...... \fontsize{13pt}{16.9pt} (size, line height) and \setstretch{1.3} (line spacing).
%   Paragraphs ..... \parindent (first-line indent) and \parskip (gap between paragraphs).
%   Heading look ... the \titleformat blocks at the end (chapter title size, "CHAPTER n:" label).
%   Page header .... the \fancyhead / \fancyfoot lines (what is printed at the top and bottom).
%   Table of contents depth ... \setcounter{tocdepth}{2}: 1 = chapters, 2 = + sections, 3 = + subsections.
%   Language names . the \renewcommand lines under "Names generated by the class".
%   Extra package .. add a \usepackage{...} line in the matching section.
%   Own shortcut ... add \newcommand{\name}{text} at the very end of this file.
% Change one thing at a time and rebuild, so you know what caused a problem.
% ======================================================================

% ---------------------------------------------------------------- Fonts
\usepackage{fontspec}
\IfFontExistsTF{Times New Roman}
  {\setmainfont{Times New Roman}}
  {\setmainfont{TeX Gyre Termes}}       % free Times clone, used if Times is missing

% ---------------------------------------------------------------- Page
\usepackage[a4paper, top=20mm, bottom=20mm, left=25mm, right=15mm]{geometry}
\usepackage{setspace}
\usepackage{indentfirst}                % indent the first paragraph after a heading too

% ---------------------------------------------------------------- Graphics
\usepackage{graphicx}
\usepackage{pdfpages}                   % insert whole PDF pages
\usepackage{adjustbox}
\usepackage[dvipsnames]{xcolor}
\usepackage{tikz}
\usetikzlibrary{calc, positioning, arrows.meta, fit, backgrounds, shapes.geometric, decorations.pathmorphing}

% A missing image shows a framed box with its file name instead of an error.
\let\origincludegraphics\includegraphics
\renewcommand{\includegraphics}[2][]{%
  \IfFileExists{#2}
    {\origincludegraphics[#1]{#2}}
    {\fbox{\parbox{0.8\linewidth}{\centering\small\ttfamily missing: \detokenize{#2}}}}}

% ---------------------------------------------------------------- Math, units, text helpers
\usepackage{amsmath}
\usepackage{amssymb}
\usepackage{siunitx}
\usepackage{csquotes}
\usepackage{enumitem}
\usepackage{etoolbox}
\usepackage{changepage}                 % adjustwidth environment

% ---------------------------------------------------------------- Tables and floats
\usepackage{booktabs}
\usepackage{tabularx}
\usepackage{longtable}
\usepackage{float}                      % [H] placement
\usepackage{caption}
\usepackage{subcaption}

% ---------------------------------------------------------------- Links
\usepackage{url}
\usepackage[colorlinks=true]{hyperref}
\hypersetup{                       % all values come from info.txt
  pdftitle  = {\thesistitle},
  pdfauthor = {\authorname},
  linkcolor = \hyperlinkcolor,         % internal links: contents, \ref
  citecolor = \hyperlinkcolor,         % citations [1]
  urlcolor  = \hyperlinkcolor,         % \url{...}
  filecolor = \hyperlinkcolor}

% ---------------------------------------------------------------- Source code
\usepackage{listings}
\lstset{
    basicstyle=\ttfamily\small,
    breaklines=true,
    backgroundcolor=\color{gray!10},
    keywordstyle=\color{blue},
    commentstyle=\color{gray},
    stringstyle=\color{red},
    numbers=left,
    numberstyle=\tiny\color{gray},
    frame=single
}
\lstnewenvironment{pythoncode}{\lstset{language=Python}}{}   % usage: \begin{pythoncode} ... \end{pythoncode}

% ---------------------------------------------------------------- Names generated by the class
\renewcommand{\contentsname}{TABLE OF CONTENTS}
\renewcommand{\bibname}{REFERENCES}
\renewcommand{\listfigurename}{LIST OF FIGURES}
\renewcommand{\listtablename}{LIST OF TABLES}
\renewcommand{\figurename}{Figure}
\renewcommand{\tablename}{Table}
\renewcommand{\appendixname}{Appendix}

% ---------------------------------------------------------------- Text layout
\setlength{\parindent}{1.5em}
\setlength{\parskip}{3pt}
\fontsize{13pt}{16.9pt}\selectfont
\setstretch{1.3}

% ---------------------------------------------------------------- Contents and lists
\usepackage{titletoc}
\usepackage{tocloft}
\setcounter{tocdepth}{2}                % 2 = chapter > section > subsection
\setcounter{secnumdepth}{2}

\captionsetup{labelsep=colon}           % "Figure 2.1: Title"
\renewcommand{\cftfigpresnum}{Figure~}
\renewcommand{\cftfigaftersnum}{:}
\setlength{\cftfignumwidth}{2.4cm}
\renewcommand{\cfttabpresnum}{Table~}
\renewcommand{\cfttabaftersnum}{:}
\setlength{\cfttabnumwidth}{2.4cm}
\renewcommand{\cftchappresnum}{CHAPTER~}
\renewcommand{\cftchapaftersnum}{:}
\setlength{\cftchapnumwidth}{3.0cm}

% ---------------------------------------------------------------- Header and footer
\usepackage{fancyhdr}
\pagestyle{fancy}
\fancyhf{}
\fancyhead[L]{\nouppercase{\leftmark}}  % "CHAPTER n"
\fancyfoot[C]{\thepage}
\renewcommand{\headrulewidth}{0pt}
\setlength{\headheight}{17pt}
\renewcommand{\chaptermark}[1]{\markboth{CHAPTER \thechapter}{}}

% ---------------------------------------------------------------- Headings
\usepackage{titlesec}
\titleformat{\chapter}[block]           % "CHAPTER 1: TITLE", centered, uppercase
  {\centering\fontsize{16pt}{18pt}\bfseries}
  {CHAPTER \thechapter:}{0.5em}{\MakeUppercase}
\titlespacing*{\chapter}{0pt}{0pt}{1cm}

\titleformat{\paragraph}[block]         % \paragraph{Title:} is its own line
  {\bfseries}{}{0pt}{}
\titlespacing*{\paragraph}{0pt}{0.6\baselineskip}{0.15\baselineskip}
